
% !TEX root = ../main.tex

% Local Variables:
% TeX-master: "../main"
% End:
% chktex-file 26

%%%%%%%%%%%%%%%%%%%%%%%%%%%%%%% Header %%%%%%%%%%%%%%%%%%%%%%%%%%%%%%%%%%%%%%%%%%%%
\begin{minipage}[l]{0.42\textwidth}
    \includegraphics[width=1\textwidth]{img/logo-UNAMBA.png}
\end{minipage}
\hfill
\begin{minipage}[c]{0.5\textwidth}
    \begin{flushright}
	\large{\textbf{Unidad \#1}}\\
	\large{Lectures on Física III}\\
	\large{06 de octubre del 2025. Haquira, Apurimac}\\
        % \large{\textbf{Student:} Huallpa Aimituma Josué David}
    \end{flushright}
\end{minipage}
\begin{tikzpicture}
    \draw[gray,thick] (-5.5,0)--(14,0);
\end{tikzpicture}


 %%%%%%%%%%%%%%%%%%%%%%%% INICIO DEL CONTENIDO EN DOS COLUMNAS %%%%%%%%%%%%%%%%%%%%%
  
 \begin{multicols}{2}

    \begin{excercise}[][][]{ex:k29}{ \textbf{(Vector Coord. Cilindricas)}\\
            Encuentre una expresión para el vector unitario del vector $\vec{A}$ mostrado en la figura en coordenadas cilíndricas.
            \begin{figure}[H]
                \centering
                \includegraphics[width=0.7\linewidth]{img/01_electric-field/1.png}
            \end{figure}
         }
    \end{excercise}

    % \begin{excercise}[][][$S=\displaystyle{\frac{5\pi}{2}}$]{ex:k29}{ \textbf{(Área cilindrica)}\\
    % %         Calcule el área de una superficie de una superficie cilindrica descrita por $r=5$, $30^\circ\leq\phi\leq 60^\circ$ y $0\leq z\leq 3$            
    % \begin{figure}[H]
    %             \centering
    %             \includegraphics[width=0.7\linewidth]{img/01_electric-field/2.png}
    %         \end{figure}
    %      }
    % \end{excercise}
    %
    % \begin{excercise}[][][$V=471.2\, \rm{cm^3}$]{ex:k29}{ 
    %         Un cilindro circular de radio $r=5$ cm  es concéntrico con el eje $z$ y se extiende entre $z=-3$ cm y $z=3$cm. Emplee la ecuación:
    %         \begin{equation*}
    %             dv=r\,dr\,d\phi\ dz
    %         \end{equation*}
    %         para determinar el volumen del cilindro. 
    %      }
    % \end{excercise}
    %
    % \begin{excercise}[][][$S=20.7\, \rm{cm^2}$]{ex:k29}{ \textbf{(Área cilindrica)}\\
    %         La franja esférica señalada en la figura es una sección de una esfera de $3$ cm de radio. Calcule el área de la franja.           
    %         \begin{figure}[H]
    %             \centering
    %             \includegraphics[width=0.8\linewidth]{img/01_electric-field/3.png}
    %         \end{figure}
    %      }
    % \end{excercise}
    %
    % \begin{excercise}[][][$P_1=P_1(5, 306.9^\circ, 3)$; $\vec{A}=3.6\vec{e}_r-0.2\vec{e}_\phi+4\vec{k}$]{ex:k29}{ 
    %         Dados el punto $P_1 (3, 4, 3)$ y el vector
    %         \begin{equation*}
    %             \vec{A}=2\vec{\imath}-3\vec{\jmath} +4\vec{k}
    %         \end{equation*}
    %         definidos en coordenadas cartesianas, exprese $P_1$ y $\vec{A}$ en coordenadas cilíndricas y evalúe $\vec{A}$ en $P_1$ .
    %      }
    % \end{excercise}
    %
    % \begin{excercise}[][][$\vec{A}=R\vec{e}_R-R\sin{\phi}\vec{e}_\theta$]{ex:k29}{ 
    %         Exprese el vector 
    %         \begin{equation*}
    %             \vec{A}=(x+y)\vec{\imath}-(y-x)\vec{\jmath} +z\vec{k}
    %         \end{equation*}
    %         en coordenadas esféricas.
    %     }
    % \end{excercise}
    %
    % \begin{excercise}[][][$P=P(4, 2\pi/3, \pi/3)$]{ex:k29}{ 
    %         El punto $P(2\sqrt{3},\pi/3, -2)$ se daen coordenadas cilíndricas. Exprese $P$ en coordenadas esféricas.        
    %     }
    % \end{excercise}
    %
    % \begin{excercise}[][][$\vec{A}=R\vec{e}_R-R\sin{\phi}\vec{e}_\theta$]{ex:k29}{ 
    %         Trasnforme el vector 
    %         \begin{equation*}
    %             \vec{A}=(x+y)\vec{\imath}-(y-x)\vec{\jmath} +z\vec{k}
    %         \end{equation*}
    %         de coordenadas cartesianas a cilíndricas.
    %     }
    % \end{excercise}
\end{multicols}
